\documentclass[11pt,a4paper]{article}
\usepackage[utf8]{inputenc}
\usepackage[T1]{fontenc}
\usepackage[ngerman]{babel}
\usepackage{geometry}
\usepackage{array}
\usepackage{tabularx}
\usepackage{booktabs}
\usepackage{amsmath}
\usepackage{amssymb}
\usepackage{graphicx}
\usepackage{tikz}
\usepackage{fancyhdr}
\usepackage{lastpage}
\usepackage{helvet}
\renewcommand{\familydefault}{\sfdefault}

\geometry{left=2cm, right=2cm, top=2.5cm, bottom=2cm}

\pagestyle{fancy}
\fancyhf{}
\fancyhead[L]{Stundenleistung Physik -- Kernphysik}
\fancyhead[R]{Klasse 10}
\fancyfoot[C]{Seite \thepage\ von \pageref{LastPage}}

\newcommand{\punkte}[1]{\hfill\textbf{/#1 BE}}
\newcommand{\antwortzeile}{\vspace{0.5cm}\hrulefill\par}
\newcommand{\antwortfeld}[1]{\vspace{#1}\hrulefill\par\vspace{0.2cm}\hrulefill\par\vspace{0.2cm}\hrulefill\par}

\begin{document}

% === DECKBLATT ===
\begin{center}
\Large\textbf{STUNDENLEISTUNG PHYSIK}\\[0.3cm]
\large Kernphysik -- Klasse 10\\[1cm]
\end{center}

\begin{tabular}{@{}p{8cm}p{6cm}@{}}
Name: \hrulefill & Datum: \hrulefill \\[0.5cm]
Klasse: \hrulefill & \\
\end{tabular}

\vspace{0.5cm}

\begin{tabular}{|p{4cm}|p{4cm}|p{4cm}|}
\hline
\textbf{Bearbeitungszeit} & \textbf{Hilfsmittel} & \textbf{Punkte} \\
\hline
25 Minuten & Taschenrechner, PSE & \hspace{1cm}/29 BE \\
\hline
\end{tabular}

\vspace{0.5cm}

\textbf{Notenschlüssel:}
\begin{center}
\begin{tabular}{|c|c|c|c|c|c|c|}
\hline
\textbf{Note} & 1 & 2 & 3 & 4 & 5 & 6 \\
\hline
\textbf{BE} & 28--29 & 23--27 & 17--22 & 12--16 & 6--11 & 0--5 \\
\hline
\textbf{\%} & $\geq$96 & $\geq$79 & $\geq$59 & $\geq$41 & $\geq$21 & $<$21 \\
\hline
\end{tabular}
\end{center}

\vspace{0.3cm}
\hrule
\vspace{0.5cm}

% === AUFGABE 1 ===
\section*{Aufgabe 1: Atommodelle \punkte{4}}

\textbf{a)} Nenne die vier historischen Atommodelle in chronologischer Reihenfolge. \punkte{2}

\begin{tabular}{@{}p{3.5cm}p{3.5cm}p{3.5cm}p{3.5cm}@{}}
1. \hrulefill & 2. \hrulefill & 3. \hrulefill & 4. \hrulefill \\
\end{tabular}

\vspace{0.5cm}

\textbf{b)} Welches Experiment führte zur Entwicklung des Kern-Hülle-Modells? \punkte{1}

\antwortzeile

\vspace{0.3cm}

\textbf{c)} Nenne die wichtigste Erkenntnis aus diesem Experiment. \punkte{1}

\antwortzeile

\vspace{0.5cm}

% === AUFGABE 2 ===
\section*{Aufgabe 2: Strahlungsarten \punkte{6}}

\textbf{a)} Fülle die Tabelle aus: \punkte{3}

\begin{center}
\begin{tabular}{|p{3cm}|p{2.5cm}|p{2.5cm}|p{2.5cm}|}
\hline
\textbf{Eigenschaft} & $\alpha$ & $\beta$ & $\gamma$ \\
\hline
Ladung & & & \\[0.8cm]
\hline
Abschirmung durch & & & \\[0.8cm]
\hline
\end{tabular}
\end{center}

\vspace{0.3cm}

\textbf{b)} Stelle die Zerfallsgleichung auf: ${}^{226}_{88}$Ra zerfällt unter $\alpha$-Strahlung. \punkte{2}

\vspace{1cm}
\begin{center}
${}^{226}_{88}\text{Ra} \longrightarrow$ \hrulefill $+$ \hrulefill
\end{center}

\vspace{0.3cm}

\textbf{c)} Erkläre, warum $\alpha$-Strahlung bei innerer Bestrahlung besonders gefährlich ist. \punkte{1}

\antwortzeile

\vspace{0.5cm}

% === AUFGABE 3 ===
\section*{Aufgabe 3: Halbwertszeit \punkte{6}}

Ein radioaktives Präparat hat zu Beginn 800 radioaktive Kerne. Die Halbwertszeit beträgt 4 Stunden.

\vspace{0.3cm}

\textbf{a)} Wie viele Kerne sind nach 12 Stunden noch vorhanden? \punkte{3}

Rechnung:
\antwortfeld{1.5cm}

Ergebnis: \hrulefill\ Kerne

\vspace{0.5cm}

\textbf{b)} Nach wie vielen Stunden sind noch 50 Kerne vorhanden? \punkte{3}

Rechnung:
\antwortfeld{1.5cm}

Ergebnis: \hrulefill\ Stunden

\vspace{0.5cm}

% === AUFGABE 4 ===
\section*{Aufgabe 4: Kernreaktionen \punkte{5}}

\textbf{a)} Erkläre den Begriff ,,Kettenreaktion`` bei der Kernspaltung. \punkte{2}

\antwortfeld{1cm}

\vspace{0.3cm}

\textbf{b)} Nenne zwei Unterschiede zwischen Kernspaltung und Kernfusion. \punkte{2}

1. \hrulefill

\vspace{0.3cm}

2. \hrulefill

\vspace{0.5cm}

\textbf{c)} Wo findet Kernfusion auf natürliche Weise statt? \punkte{1}

\antwortzeile

\vspace{0.5cm}

% === AUFGABE 5 ===
\section*{Aufgabe 5: Strahlenschutz \punkte{4}}

\textbf{a)} Nenne die drei Schutzmaßnahmen gegen ionisierende Strahlung (die ,,3 A's``). \punkte{1,5}

\begin{tabular}{@{}p{4cm}p{4cm}p{4cm}@{}}
1. \hrulefill & 2. \hrulefill & 3. \hrulefill \\
\end{tabular}

\vspace{0.5cm}

\textbf{b)} Ein Arzt arbeitet mit $\gamma$-Strahlern. Erkläre, wie er sich schützen kann. \punkte{2,5}

\antwortfeld{1cm}

\vspace{0.5cm}

% === AUFGABE 6 ===
\section*{Aufgabe 6: Bewertung \punkte{4}}

Ein Patient soll wegen Rückenschmerzen geröntgt werden. Der Arzt schlägt alternativ ein CT vor, das genauere Bilder liefert, aber eine höhere Strahlendosis bedeutet.

\vspace{0.3cm}

\textbf{a)} Nenne je einen Vorteil von Röntgen und CT. \punkte{2}

Röntgen: \hrulefill

\vspace{0.3cm}

CT: \hrulefill

\vspace{0.5cm}

\textbf{b)} Nimm Stellung: Sollte der Patient das CT machen lassen? Begründe mit Bezug auf Nutzen und Risiko. \punkte{2}

\antwortfeld{1.5cm}

\vspace{1cm}

% === ZUSATZAUFGABE ===
\hrule
\vspace{0.3cm}

\section*{Zusatzaufgabe (Bonus, 5 BE)}

\textit{Nur bearbeiten, wenn du vor Ablauf der Zeit fertig bist. Zählt nicht negativ!}

\vspace{0.3cm}

In einem Kernreaktor werden pro Sekunde $3 \cdot 10^{19}$ U-235-Kerne gespalten.
Pro Spaltung werden 200 MeV frei.

\textbf{Daten:} 1 MeV = $1{,}6 \cdot 10^{-13}$ J

\vspace{0.3cm}

\textbf{a)} Berechne die thermische Leistung des Reaktors in Watt. \punkte{3}

\antwortfeld{2cm}

\textbf{b)} Erkläre, warum die elektrische Leistung geringer ist. \punkte{2}

\antwortfeld{1cm}

\vspace{0.5cm}

\begin{center}
\textit{--- Ende der Stundenleistung ---}
\end{center}

\newpage

% === LÖSUNGSBLATT ===
\section*{Lösungen (für Lehrkraft)}

\subsection*{Aufgabe 1 (4 BE)}
\begin{itemize}
    \item a) Dalton, Thomson, Rutherford, Bohr (je 0,5 BE)
    \item b) Rutherford-Streuversuch (1 BE)
    \item c) Das Atom hat einen kleinen, positiven Kern / ist größtenteils leer (1 BE)
\end{itemize}

\subsection*{Aufgabe 2 (6 BE)}
\begin{itemize}
    \item a) $\alpha$: +2, Papier/Haut | $\beta$: --1, Aluminium | $\gamma$: 0, Blei/Beton (je 0,5 BE)
    \item b) ${}^{226}_{88}\text{Ra} \rightarrow {}^{222}_{86}\text{Rn} + {}^{4}_{2}\text{He}$ (2 BE)
    \item c) Hohes Ionisierungsvermögen $\rightarrow$ große Schäden im Gewebe (1 BE)
\end{itemize}

\subsection*{Aufgabe 3 (6 BE)}
\begin{itemize}
    \item a) 12h $\div$ 4h = 3 HWZ $\rightarrow$ 800 $\cdot$ $(1/2)^3$ = 800 $\cdot$ 1/8 = \textbf{100 Kerne} (3 BE)
    \item b) 800 $\rightarrow$ 400 $\rightarrow$ 200 $\rightarrow$ 100 $\rightarrow$ 50 = 4 HWZ $\rightarrow$ 4 $\cdot$ 4h = \textbf{16 Stunden} (3 BE)
\end{itemize}

\subsection*{Aufgabe 4 (5 BE)}
\begin{itemize}
    \item a) Bei der Spaltung werden Neutronen frei, die weitere Kerne spalten. So vermehrt sich die Anzahl exponentiell. (2 BE)
    \item b) Z.B.: Spaltung: schwere $\rightarrow$ leichtere Kerne, Fusion: leichte $\rightarrow$ schwerere / Atommüll vs. kaum Abfall / ausgereift vs. Forschung (je 1 BE)
    \item c) In der Sonne / in Sternen (1 BE)
\end{itemize}

\subsection*{Aufgabe 5 (4 BE)}
\begin{itemize}
    \item a) Abstand, Abschirmung, Aufenthaltszeit (je 0,5 BE)
    \item b) Großer Abstand (Zange), Bleischürze/-wand, Aufenthalt minimieren, Dosimeter (2,5 BE)
\end{itemize}

\subsection*{Aufgabe 6 (4 BE)}
\begin{itemize}
    \item a) Röntgen: geringere Strahlendosis, schneller, günstiger / CT: genauere 3D-Bilder, bessere Diagnose (je 1 BE)
    \item b) Begründete Stellungnahme mit Nutzen-Risiko-Abwägung (2 BE)
\end{itemize}

\subsection*{Zusatzaufgabe (5 BE Bonus)}
\begin{itemize}
    \item a) $P = 3 \cdot 10^{19} \cdot 200 \cdot 1{,}6 \cdot 10^{-13}$ J/s = $9{,}6 \cdot 10^{8}$ W $\approx$ 960 MW (3 BE)
    \item b) Wirkungsgrad $<$ 100\%: Wärme geht verloren (Kühlturm), Turbine/Generator nicht 100\% effizient (2 BE)
\end{itemize}

\end{document}
